\section{Solving the allocation, and certifying it}
\label{sec:algorithm}

Problem~\eqref{eq:allocation} is solved twice: once from above, by a descent
that returns an attainable schedule, and once from below, by a relaxation that
bounds what any schedule could achieve. The pair is what licenses the language
used for the result.

\subsection{Attainable schedule: budgeted coordinate descent}
\label{sec:algorithm-descent}

Introduce the multiplier $\lambda$ on the budget and minimize
$J(\mathbf u)+\lambda\sum_q W_q N_q^2$ by sweeping the \emph{decision
intervals}. Holding all other intervals fixed, the terms involving $N_q$ are

\begin{equation}
  N_q \;\leftarrow\; \argmin_{N\in\mathcal{N}}\;
  \Bigl[\;
    \underbrace{\sum_{i,k\in I_q}\mathbf u_i(N)^{\top}\mathbf Q_{ik}\,
      \mathbf u_k(N)}_{\text{within block, Eq.~\eqref{eq:within-block}}}
    \;+\; 2\sum_{i\in I_q}\mathbf u_i(N)^{\top}\mathbf c_i^{(-q)}
    \;+\; \lambda\,W_q\,N^2
  \;\Bigr] ,
  \label{eq:coordinate-update}
\end{equation}

with the interval time weight $W_q$ of Eq.~\eqref{eq:interval-weight} in place
of the sample count, and
$\mathbf c_i^{(-q)}=\mathbf c_i-\sum_{k\in I_q}\mathbf Q_{ik}\mathbf u_k$
from Eq.~\eqref{eq:cross-term}. The first group of terms is the block's own
contribution --- a magnitude, which is what a separable method would see ---
and the second is the coupling: it rewards a defect pointing \emph{against} the
accumulated displacement. That second term is the whole difference between this
and a sensitivity-weighted allocation.

\begin{algorithm}[!htbp]
\caption{Budgeted coordinate descent for Eq.~\eqref{eq:allocation}}
\label{alg:descent}
\begin{algorithmic}[1]
\Require tabulated $\mathbf u_i(N)$ for $i=1..M$, $N\in\mathcal{N}$; suffix
  blocks $\mathbf A_i$; grouping $g$; budget $B$; start set $\mathcal{S}$
  (Table~\ref{tab:starts})
\State $\lambda \gets$ bracket from the separable solution
\Repeat \Comment{outer: bisection on the budget}
  \For{each start $s\in\mathcal{S}$} \Comment{multi-start, \emph{not} one start}
    \State $N \gets s$
    \Repeat \Comment{inner: sweeps}
      \State refresh prefix sums of $\mathbf u$ and suffix sums of
        $\mathbf A_k\mathbf u_k$
      \For{$q = 1 \dots K_{\mathrm{dec}}$}
        \State update $N_q$ by Eq.~\eqref{eq:coordinate-update}; patch the two
          running sums over $I_q$
      \EndFor
    \Until{no interval changes, or sweep limit}
    \State record $(J(N),\,N)$
  \EndFor
  \State $N^\star(\lambda) \gets \argmin_s J$ \Comment{best of starts, at this
    $\lambda$}
  \State $\lambda \gets$ bisect on $\sum_q W_q N_q^{\star2} - B$
\Until{budget met to tolerance \ph{value}}
\State \Return $N^\star$, $J(N^\star)$, and the spread of $J$ across
  $\mathcal{S}$
\end{algorithmic}
\end{algorithm}

The multi-start loop sits \emph{inside} the bisection rather than outside it:
each $\lambda$ is solved from scratch from all starts. That is more expensive
than warm-starting from the previous $\lambda$, and it is deliberate --- a
warm start would make the returned schedule depend on the bisection path, which
is precisely the failure mode the monotonicity check below is meant to detect.

Each sweep is $\mathcal{O}(M\lvert\mathcal{N}\rvert)$ by
Eqs.~\eqref{eq:cross-term} and~\eqref{eq:within-block} --- grouping changes the
number of decision variables but not the sweep cost --- and no field evaluation
occurs inside the loop: the $\mathbf u_i(N)$ table is built once per orbit.

\paragraph{Summation.} The prefix sums are accumulated in compensated
(Neumaier) arithmetic. This is not routine caution: $\mathbf S_j$ is a sum of
order $10^5$ signed terms whose cancellation is the very quantity the method
exploits, so the summation is ill-conditioned by construction, and the ratio of
$\sum_i\lVert\mathbf u_i\rVert$ to $\lVert\mathbf S_M\rVert$ is reported per
orbit as the conditioning diagnostic. \wpref{WP1}

The objective is not convex in the integer variables, so the descent is
started from several points and the spread across starts is reported rather
than suppressed \wpref{WP3}.

\paragraph{The outer bisection needs a monotonicity check.} For a global
minimizer the attained work $W(\lambda)$ is non-increasing in $\lambda$, which
is what licenses bisecting on it. What Algorithm~\ref{alg:descent} returns is a
best-of-starts local solution, and nothing guarantees that the basin it lands
in varies monotonically. We therefore sweep $\lambda$ on a dense logarithmic
grid before any campaign run and verify that the best-of-starts $W(\lambda)$ is
non-increasing. Where it is, bisection is used. Where it is not, the bisection
is abandoned for that orbit and the schedule is chosen as the feasible grid
point with the smallest $J$, which is a weaker procedure but not an unsound
one; the orbits on which this happens are counted and reported. \wpref{WP3}

\begin{table}[!htbp]
  \centering\small
  \caption{Starting schedules for Algorithm~\ref{alg:descent} and the spread of
  the attained objective across them. Declared before the runs.}
  \label{tab:starts}
  \begin{tabular}{@{}llc@{}}
    \toprule
    Start & Rationale & Median $J/J_{\mathrm{best}}$ \\
    \midrule
    \Fop    & constant degree at the budget & \ph{value} \\
    \Rrad   & budget-calibrated radial endpoint & \ph{value} \\
    \Rint   & interior member & \ph{value} \\
    \Asens  & separable sensitivity-weighted solution & \ph{value} \\
    random $\times\,8$ & spread & \ph{value}--\ph{value} \\
    \bottomrule
  \end{tabular}
\end{table}

\subsection{Certificate: convex-hull relaxation and a Frank--Wolfe gap}
\label{sec:algorithm-certificate}

Calling the descent result a benchmark requires knowing how far it can be from
the best possible schedule. Relax each \emph{decision interval} to the convex
hull of its candidate contributions,

\begin{equation}
  \mathbf u_i \;=\; \sum_{N\in\mathcal{N}}\theta_{g(i)N}\,\mathbf u_i(N),
  \qquad
  \theta_{qN}\ge0, \quad \sum_N \theta_{qN}=1,
  \qquad
  \sum_q W_q \sum_N \theta_{qN}N^2 \;\le\; B .
  \label{eq:relaxation}
\end{equation}

Every integer schedule is feasible here with $\theta$ an indicator and attains
the same objective, so the relaxed optimum $J^{*}_{\mathrm{rel}}$ is a lower
bound on the integer optimum $J^{*}$. The relaxed problem minimizes the convex
quadratic $J$ over a convex set, and its linear subproblem

\begin{equation}
  \mathbf s \;\in\; \argmin_{\theta \text{ feasible}}\;
  \langle \nabla J(\mathbf u), \mathbf u(\theta)\rangle
  \label{eq:fw-subproblem}
\end{equation}

decouples across decision intervals under a single multiplier on the budget. It
therefore has the same epoch-separable, single-multiplier \emph{structure} as
the sensitivity-weighted allocation of Section~\ref{sec:benchmark-sens}, but it
is not the same problem: its per-interval objective is
$\langle\nabla_q J,\mathbf u_q(N)\rangle$ rather than
$\sum_{i\in I_q}\Delta t_i\Delta\mathbf a_i^{\top}\mathbf K_i\Delta\mathbf a_i$,
and
$\nabla J = 2\mathbf Q\mathbf u$ depends on the current relaxed iterate.
Frank--Wolfe therefore applies with an
$\mathcal{O}(M\lvert\mathcal{N}\rvert)$ iteration
\citep{frank1956algorithm,jaggi2013revisiting}, and its duality gap supplies
the bound directly:

\begin{equation}
  J(\mathbf u) + \langle\nabla J(\mathbf u),\,\mathbf s-\mathbf u\rangle
  \;\le\; J^{*}_{\mathrm{rel}} \;\le\; J^{*}
  \;\le\; J(\mathbf u_{\mathrm{descent}}) .
  \label{eq:certificate}
\end{equation}

Write $L_{\mathrm{FW}}$ for the best (largest) value the left-hand side attains
over the iteration budget and $J_{\mathrm{desc}}$ for the descent objective.
Because the objective is a squared error, a gap in $J$ and a gap in $E$ are not
the same number, so both are defined and the threshold is fixed on the second:

\begin{equation}
  g_J \;=\; \frac{J_{\mathrm{desc}}-L_{\mathrm{FW}}}{J_{\mathrm{desc}}},
  \qquad
  g_E \;=\; 1-\sqrt{\frac{L_{\mathrm{FW}}}{J_{\mathrm{desc}}}}
  \;=\; 1-\sqrt{1-g_J} .
  \label{eq:certified-gap}
\end{equation}

$g_E$ is the quantity the paper's claims are in --- it states that the
descent schedule's position error is within $g_E$ of the best error any
schedule of this class could achieve --- and the naming rule below is bound to
it. A $10\%$ error gap corresponds to $g_J\approx19\%$.

Two degenerate cases are handled by rule rather than by judgement. The
Frank--Wolfe bound is valid but can be vacuous early in the iteration, where
the left-hand side of Eq.~\eqref{eq:certificate} may fall below zero; since
$J\ge0$ always, $L_{\mathrm{FW}}$ is taken as
$\max\{0,\ \text{best bound}\}$. If $L_{\mathrm{FW}}=0$ when the iteration
budget is exhausted, the certificate is vacuous, no gap is reported for that
orbit, and the orbit is counted in a separate column rather than dropped.

\paragraph{The subproblem must be solved exactly.} The bound in
Eq.~\eqref{eq:certificate} is valid only if $\mathbf s$ attains the minimum of
Eq.~\eqref{eq:fw-subproblem} over the \emph{relaxed} feasible set. Stopping the
multiplier bisection at the nearest schedule whose integer budget falls below
$B$ does not do that: it returns a feasible point, not a minimizer, and the
resulting quantity is no longer a lower bound. Because
Eq.~\eqref{eq:relaxation} permits fractional mixing, the budget can be met with
equality --- at the critical multiplier the solution mixes the two adjacent
degrees of the affected intervals with the weights that saturate the
constraint. The implementation constructs that mixture explicitly and asserts
$\lvert\sum_q W_q\sum_N\theta_{qN}N^2 - B\rvert$ at machine
precision on every iteration; a run in which the assertion fails produces no
certificate.

\paragraph{What the certificate does and does not license.} It bounds the
distance to the best schedule \emph{of this problem}: the first-order
displacement along a fixed reference trajectory, on a fixed accumulation grid,
over a discrete degree set. It says nothing about the nonlinear trajectory,
which Section~\ref{sec:campaign-fixedpoint} tests separately. We adopt the
naming rule fixed before the runs: the quantity is called an \emph{oracle} only
if the median $g_E$ of Eq.~\eqref{eq:certified-gap} is below $0.10$ and no more
than \ph{n} orbits return a vacuous certificate; otherwise it is reported as a
\emph{linearized trajectory-aware allocation benchmark} and every ratio taken
against it is read as conservative, in the same sense the previous paper's
force-level Lagrangian benchmark was.

\dnote{One of the two names has to be chosen once Table~\ref{tab:certificate}
exists, and the choice is bound to the gap number rather than to the result.
If the gap is large, delete the word ``oracle'' from the whole manuscript ---
including the abstract --- rather than qualifying it in place.}

\begin{table}[!htbp]
  \centering\small
  \caption{Certified gap between the Frank--Wolfe lower bound and the
  coordinate-descent schedule, by design and budget. \ph{fill from WP7}}
  \label{tab:certificate}
  \begin{tabular}{@{}lccccccc@{}}
    \toprule
    Design & $\beta$ & $L_{\mathrm{FW}}$ & $J_{\mathrm{desc}}$ & $g_J$ & $g_E$
      & worst $g_E$ & vacuous \\
    \midrule
    A & $1.00$ & \ph{value} & \ph{value} & \ph{pct} & \ph{pct} & \ph{pct} & \ph{n} \\
    B & $1.00$ & \ph{value} & \ph{value} & \ph{pct} & \ph{pct} & \ph{pct} & \ph{n} \\
    \bottomrule
  \end{tabular}
\end{table}

\subsection{Cost of the solve}

\ph{one paragraph} Table build per orbit; sweep count to convergence;
Frank--Wolfe iterations to a useful gap; memory of the vector table at the
converged accumulation grid. None of this is charged to the flight budget ---
the benchmark is not flown --- but it decides what is feasible to run across a
campaign, and the controller of Section~\ref{sec:controller} inherits a
reduced form of it, which \emph{is} charged. \wpref{WP1, WP3, WP7}
